\chapter{Metody poprawy jakości na małym zbiorze danych}

W niniejszym rozdziale zaprezentowano przegląd literatury dotyczącej metod radzenia sobie z problemem niedoboru danych treningowych w głębokich sieciach neuronowych. W pierwszej kolejności zdefiniowano wyzwania związane z uczeniem modeli w reżimie małego zbioru danych, odnosząc standardy literaturowe do specyfiki problemu badanego w niniejszej pracy. Następnie omówiono dwie klasy technik będące przedmiotem eksperymentów: uczenie pół-nadzorowane z wykorzystaniem pseudo-etykietowania oraz syntetyczną augmentację danych z użyciem modeli dyfuzyjnych. Rozdział zamyka analiza najnowszych prac badawczych (ang.\ \emph{State-of-the-Art}), wykorzystujących wspomniane metody w zadaniach widzenia komputerowego.

\section{Wyzwania reżimu ,,małego zbioru danych''}

Skuteczność głębokich sieci neuronowych zwykle rośnie wraz ze zwiększaniem liczby danych treningowych. Zjawisko to, opisywane w literaturze jako ,,nierozsądna skuteczność danych'' (ang.\ \emph{unreasonable effectiveness of data}) \cite{Sun2017}, powoduje, że współczesne modele wizyjne wymagają bardzo dużych zbiorów treningowych. Aby zdefiniować pojęcie ,,małego zbioru'', należy odnieść się do standardów branżowych. Współczesne modele wizyjne trenowane są wstępnie na dużych zbiorach danych takich jak ImageNet, zawierającym ponad 14 milionów oznaczonych obrazów \cite{Deng2009}, czy MS COCO, składającym się z ponad 330 tysięcy zdjęć i 1.5 miliona instancji obiektów \cite{Lin2014}. 

W literaturze taki problem określa się jako reżim ograniczonej liczby danych (ang.\ \emph{low-data regime}). Obejmuje on skrajne przypadki uczenia nielicznych próbek (ang.\ \emph{few-shot learning}), gdzie model dysponuje zaledwie od 1 do 50 przykładów na klasę \cite{Wang2020FewShot}, a także szerszą kategorię małych zbiorów, liczących od kilkuset do kilku tysięcy obrazów. W kontekście modeli wizyjnych posiadających dużą liczbę parametrów (takich jak głębokie warianty YOLO czy modele typu transformer), nawet zbiór liczący kilka tysięcy oznaczonych próbek często uznawany jest za niewystarczający do optymalizacji sieci od podstaw.

W praktycznych zastosowaniach inżynierskich i przemysłowych pozyskanie danych na skalę zbioru COCO jest najczęściej niemożliwe. Problem badawczy zdefiniowany w tej pracy stanowi klasyczny przykład operowania w reżimie małego zbioru danych. Zadanie polega na detekcji specyficznych struktur ochronnych (worków z piaskiem) na podstawie zdjęć lotniczych pozyskiwanych przez bezzałogowe statki powietrzne (UAV). Oryginalny zbiór treningowy wykorzystany w badaniach składa się zaledwie z około 210 ręcznie oznaczonych fragmentów obrazów. Niewielki rozmiar tego zbioru wynika z trudności związanych z pozyskiwaniem danych: monitorowane struktury pojawiają się niemal wyłącznie w rzadkich sytuacjach kryzysowych (zagrożenie powodziowe), a ich gęste, ręczne etykietowanie przez ekspertów jest procesem czasochłonnym i kosztownym.

Problem utrudnia również charakterystyka wizualna analizowanych obiektów. Worki z piaskiem występują z reguły w ogromnych zagęszczeniach, tworząc tzw.\ sceny gęsto upakowane (ang.\ \emph{crowded scenes}). Jak zauważają Gao i in. w swoim przekrojowym badaniu \cite{Gao2020}, detekcja obiektów w takich warunkach napotyka na istotne trudności związane z silną okluzją (obiekty zasłaniają się wzajemnie), dużą zmiennością skali wynikającą z perspektywy oraz złożonym tłem (ang.\ \emph{complex background}), które często zlewa się kolorystycznie z obiektami docelowymi.

Głównym problemem przy tak małym zbiorze danych jest przeuczenie modelu (ang.\ \emph{overfitting}). Współczesne sieci neuronowe posiadają bardzo dużą liczbę parametrów, przez co mogą łatwo dopasować się do danych treningowych. W efekcie model zamiast uczyć się ogólnych cech obiektu zaczyna zapamiętywać elementy charakterystyczne dla konkretnych próbek, takie jak tło, szum obrazu czy artefakty kamery. Taki model osiąga dobre wyniki na danych treningowych, ale gorzej radzi sobie z nowymi obrazami. Problem ten jest szczególnie widoczny przy zmianie warunków oświetleniowych lub rodzaju podłoża, gdzie mogą pojawiać się błędne detekcje.

\section{Uczenie pół-nadzorowane i pseudo-etykietowanie}

W wielu scenariuszach operacyjnych pozyskanie surowych danych (nieoznakowanych obrazów) z sensorów wizyjnych jest tanie i proste, natomiast proces ich manualnej adnotacji przez ekspertów jest kosztowny i czasochłonny. Odpowiedzią na ten problem jest uczenie pół-nadzorowane (ang.\ \emph{Semi-Supervised Learning}, SSL), które łączy w procesie optymalizacji niewielki zbiór danych oznaczonych $D_L$ (ang.\ \emph{labeled data}) oraz potencjalnie wielokrotnie większy zbiór danych nieoznakowanych $D_U$ (ang.\ \emph{unlabeled data}).

Jednym z najczęściej stosowanych paradygmatów SSL w widzeniu komputerowym jest uczenie samonadzorowane wykorzystujące pseudo-etykietowanie (ang.\ \emph{self-training with pseudo-labeling}). Klasyczny proces realizowany jest w układzie Nauczyciel-Uczeń (ang.\ \emph{Teacher-Student}). W pierwszej fazie inicjującej model Nauczyciela trenowany jest w sposób w pełni nadzorowany, bazując wyłącznie na ograniczonej puli danych z etykietami ($D_L$). Następnie Nauczyciel wykorzystywany jest do wygenerowania predykcji na zbiorze nieoznakowanym ($D_U$). Wygenerowane predykcje są następnie filtrowane według wybranych kryteriów jakości. Te, które spełniają wymagania, są wykorzystywane jako pseudo-etykiety do trenowania modelu ucznia \cite{Lee2013}.

Istotnym problemem tego podejścia jest błąd potwierdzenia (ang.\ \emph{confirmation bias}) wynikający ze zjawiska tzw.\ nadmiernej pewności modelu (ang.\ \emph{overconfidence}) \cite{Guo2017}. Głębokie sieci neuronowe wykazują tendencję do generowania wysokich wartości prawdopodobieństwa wyjściowego (pewności) nawet dla całkowicie błędnych predykcji. Jeśli Nauczyciel popełni błąd klasyfikacyjny z dużą pewnością (co jest szczególnie częste w reżimie małych zbiorów na skutek przeuczenia), Uczeń będzie trenowany na błędnych etykietach, co może pogorszyć skuteczność modelu. Z tego powodu w większości prac stosuje się odrzucanie predykcji niespełniających rygorystycznych kryteriów jakościowych. Jedną z najczęściej stosowanych metod jest progowanie pewności (ang.\ \emph{Confidence Thresholding}), dopuszczające do procesu optymalizacji wyłącznie te pseudo-etykiety, których wartość funkcji prawdopodobieństwa wyjściowego przekracza z góry zdefiniowany próg empiryczny (np.\ $0.45$ lub więcej) \cite{Sohn2020}. 

Należy jednak zauważyć, że sztywne progowanie powoduje bezpowrotną utratę części potencjalnie użytecznych informacji o trudniejszych obiektach. Współczesne podejścia opisywane w literaturze, adresujące problem niedoskonałych etykiet referencyjnych w gęstych scenariuszach, próbują ten problem optymalizować bezpośrednio w funkcji straty. Przykładowo, Nguyen i in. \cite{Nguyen2022} w dwuetapowym procesie treningowym zastosowali estymację niepewności (ang.\ \emph{uncertainty estimation}). Sieć uczy się przypisywać wygenerowanym pseudo-ramkom miarę niepewności, co działa jako automatyczny regularyzator podczas uczenia -- błąd z predykcji o wysokiej niepewności jest tłumiony, zapobiegając propagacji fałszywych wzorców do modelu Ucznia.

Samo odrzucanie błędów jest jednak niewystarczające. Jeśli model Ucznia będzie trenowany na dokładnie tych samych, niezaburzonych obrazach z pseudo-etykietami co Nauczyciel, nie nauczy się bardziej ogólnych cech obiektu, a jedynie sklonuje wiedzę bazową. Jedno z rozwiązań tego problemu zaproponowano w architekturze \emph{Noisy Student} \cite{Xie2020}. Paradygmat ten zakłada asymetryczną augmentację: Nauczyciel generuje pseudo-etykiety na oryginalnych, czystych obrazach, natomiast model Ucznia jest zmuszany do poprawnego rozpoznania tych samych obiektów na obrazach poddanych silnym transformacjom augmentacyjnym. Zastosowanie zaburzeń w przestrzeni barw (HSV) oraz strukturalnych modyfikacji pokroju algorytmów \emph{MixUp} \cite{Zhang2017} czy \emph{Mosaic} wymusza na sieci Ucznia uniezależnienie się od prostych tekstur i skupienie uwagi na solidnych, niezmienniczych cechach geometrycznych obiektu docelowego. 

Współczesne techniki \emph{self-training} nie są ponadto procesami jednorazowymi, lecz działają w cyklach iteracyjnych \cite{Zoph2020}. Po zakończeniu cyklu uczenia na zaszumionych danych, model Ucznia (który dzięki silnej augmentacji cechuje się wyższą odpornością na zakłócenia i lepszą generalizacją) awansuje do roli nowego Nauczyciela. Nowy Nauczyciel dokonuje ponownej predykcji na zbiorze $D_U$, dostarczając pseudo-etykiety o wyższej jakości (z mniejszą domieszką błędu potwierdzenia). Cały proces może być powtarzany iteracyjnie. W kolejnych etapach model generuje coraz lepsze pseudo-etykiety, co stopniowo poprawia jakość uczenia.

\section{Syntetyczna augmentacja danych i modele dyfuzyjne}

Tradycyjna augmentacja danych, obejmująca m.in.\ rotacje, przesunięcia, odbicia lustrzane czy zmiany kontrastu, jest standardową techniką przeciwdziałania przeuczeniu. Jej ograniczenie polega jednak na tym, że manipuluje ona jedynie istniejącymi pikselami, nie generując nowych informacji semantycznych. Przekształcenia te nie potrafią wygenerować nowych kątów widzenia obiektu, zmiany oświetlenia sceny czy zupełnie nowego kontekstu tła.

Rozwinięciem klasycznej augmentacji jest syntetyczna augmentacja z wykorzystaniem głębokich modeli generatywnych. Po latach dominacji sieci GAN (ang.\ \emph{Generative Adversarial Networks}), jedną z dominujących metod generowania obrazów o wysokiej jakości są modele dyfuzyjne (ang.\ \emph{Denoising Diffusion Probabilistic Models}, DDPM) \cite{Ho2020}. Proces dyfuzji polega na stopniowym, opartym na łańcuchach Markowa, dodawaniu szumu Gaussa do obrazu (proces w przód), a następnie uczeniu sieci neuronowej procedury odwrotnej -- stopniowego odszumiania. Architektury z rodziny \emph{Latent Diffusion} (takie jak Stable Diffusion) \cite{Rombach2022} pozwalają na wysoce precyzyjne generowanie danych wizyjnych warunkowanych tekstem.

Jednym z rzadziej analizowanych kierunków w literaturze dotyczącej syntetycznej augmentacji, jest wykorzystanie modeli dyfuzyjnych do generowania tzw.\ próbek negatywnych (ang.\ \emph{negative samples}). W klasycznym podejściu badawczym uwaga skupia się na sztucznym powielaniu obiektów docelowych. Tymczasem w reżimie małego zbioru danych detektory obiektów wykazują silną tendencję do przeuczenia na specyficznych teksturach tła,co prowadzi do wzrostu liczby fałszywych detekcji (ang.\ \emph{False Positives}), określanych kolokwialnie jako ,,halucynacje''.

Zamiast pobierać losowe obrazy tła z ogólnodostępnych baz danych, co często wprowadza niepożądaną lukę domenową (ang.\ \emph{domain gap}), modele dyfuzyjne pozwalają na wygenerowanie fotorealistycznych, skomplikowanych scen idealnie dopasowanych do docelowego środowiska operacyjnego. Włączenie do zbioru treningowego obrazów przedstawiających wyłącznie puste tło -- którym intencjonalnie nie przypisuje się żadnych etykiet -- działa jako dodatkowa forma regularyzacji. Zmusza to architekturę sieci do aktywnego obniżania pewności siebie w obszarach pozbawionych użytecznych cech semantycznych. Badania dotyczące trudnych przykładów negatywnych (\emph{Hard Negative Mining}) \cite{Shrivastava2016} pokazują, że regularne trenowanie modelu na złożonych tłach poprawia jego odporność na fałszywe detekcje.

\section{Przegląd i analiza pokrewnych prac badawczych}

Połączenie metod detekcji obiektów z technikami radzenia sobie z ograniczoną liczbą danych stanowi aktywnie rozwijany obszar badań. Znaczące postępy poczyniono zwłaszcza w obszarze uczenia pół-nadzorowanego dla detekcji obiektów (ang.\ \emph{Semi-Supervised Object Detection}, SSOD). W pracy Liu i in.\ zaproponowano architekturę \emph{Unbiased Teacher} \cite{Liu2021}, która proponuje rozwiązanie problemu przeuczenia głowicy klasyfikacyjnej detektora (ang.\ \emph{pseudo-labeling bias}). Autorzy wykazali, że zaktualizowanie wag Nauczyciela przy użyciu średniej wykładniczej (EMA -- \emph{Exponential Moving Average}) z wag Ucznia pozwala ustabilizować proces treningu i poprawę metryki mAP nawet o kilkanaście punktów procentowych przy dostępie do zaledwie 1\%-10\% oznakowanych danych.

Równolegle, uwaga badaczy skupia się na specyfice wykrywania obiektów w dużym zagęszczeniu. Tradycyjnie, jak udokumentowali Gao i in. w przekrojowym badaniu metod detekcji tłumu \cite{Gao2020}, problem gęstych scen rozwiązywano poprzez estymację map gęstości, rezygnując z dokładnej lokalizacji przestrzennej (ramek). Powrót do precyzyjnej detekcji obiektów stał się możliwy m.in. dzięki architekturze Counting-DETR autorstwa Nguyena i in. \cite{Nguyen2022}. Wykazali oni, że wykorzystanie dwuetapowej strategii uczenia i generowania pseudo-etykiet pozwala na efektywną detekcję mnogości obiektów docelowych z użyciem architektury opartej na transformerach, nawet gdy model inicjowany jest bardzo małą liczbą danych treningowych (ang.\ \emph{few-shot learning}). Praca ta stanowi istotny punkt odniesienia dla metod wykorzystujących generowanie pseudo-etykiet.

Równolegle, w literaturze badany jest potencjał modeli dyfuzyjnych jako generatorów danych treningowych. Badania przeprowadzone przez Azizi i in.\ \cite{Azizi2023} wykazały, że fine-tuning (dostrajanie) dużych modeli tekst-na-obraz w celu syntetycznej augmentacji poprawia generalizację modeli wizyjnych (w tym detektorów) znacznie lepiej niż tradycyjne metody augmentacyjne. Z kolei w pracach eksperymentalnych z architekturami rodziny YOLO udowodniono, że wstrzykiwanie danych wygenerowanych syntetycznie do małych zbiorów treningowych poprawia skuteczność detekcji obiektów częściowo zasłoniętych lub znajdujących się w trudnych warunkach oświetleniowych.

Mimo istnienia wielu dedykowanych publikacji dla poszczególnych architektur, wciąż brakuje kompleksowych, niezależnych badań porównujących bezpośrednio opłacalność (ang.\ \emph{trade-off}) stosowania pseudo-etykietowania w zestawieniu z syntetyczną augmentacją dyfuzyjną na nowoczesnych modelach typu \emph{anchor-free} (YOLO, RT-DETR). Jednym z celów niniejszej pracy jest częściowe uzupełnienie tej luki, poprzez stworzenie spójnego środowiska eksperymentalnego, umożliwiającego bezpośrednią i weryfikowalną statystycznie analizę obu tych podejść.

\section{Podsumowanie rozdziału}

W niniejszym rozdziale zdefiniowano główne wyzwania związane z uczeniem głębokich sieci neuronowych w reżimie małego zbioru danych, ze szczególnym uwzględnieniem problemu przeuczenia oraz błędu potwierdzenia (ang.\ \emph{confirmation bias}). Przeprowadzony przegląd literatury pozwolił na wyodrębnienie dwóch kierunków badawczych uodparniających modele na zmianę domeny. Pierwszym z nich jest iteracyjne uczenie pół-nadzorowane (wykorzystujące mechanizmy asymetrycznej augmentacji, takie jak \emph{Noisy Student}), natomiast drugim -- syntetyczna augmentacja wykorzystująca generatywne modele dyfuzyjne do wprowadzania trudnych próbek negatywnych. 

Analiza dotychczasowych prac wykazała braki w kompleksowym porównaniu tych metod dla nowoczesnych, zoptymalizowanych architektur detekcji obiektów (takich jak YOLOv26 czy RT-DETR). Wnioski płynące z przeglądu stanu wiedzy stanowią podstawę do zaprojektowania eksperymentów opisanych w dalszej części pracy. \todo{Nie wprowadzałbym w podsumowaniach rozdziałów przejścia: 'co będzie w kolejnym odcinku/rozdziale'. Uwaga nie tylko do tego miejsca. Można usunąć końcówkę ' Przejście od teorii do praktyki nastąpi w kolejnym rozdziale...'. Za to warto zostawić informacje, czemu rozdział jest ważny z perspektywy całej pracy.}